\begin{thebibliography}{33}
\bibitem{r1} Paul Beame, Toniann Pitassi, and Nathan Segerlind. Lower bounds for Lovász-Schrijver systems and beyond follow from multiparty communication complexity. SIAM J. Comput., 37(3):845–869, 2007.
\bibitem{r2} Paul Beame and Søren Riis. More on the relative strength of counting principles. In Paul Beam and Samuel R. Buss, editors, Proof Complexity and Feasible Arithmetics, Proceedings of a DIMACS Workshop, New Brunswick, New Jersey, USA, April 21-24, 1996, volume 39 of DIMACS Series in Discrete Mathematics and Theoretical Computer Science, pages 13–35. DIMACS/AMS, 1996. doi:10.1090/dimacs/039/02.
\bibitem{r3} Stephen A. Cook and Robert A. Reckhow. The relative efficiency of propositional proof systems. The Journal of Symbolic Logic, 44(1):36–50, 1979.
\bibitem{r4} William Cook, Collette R. Coullard, and Gy. Turán. On the complexity of cutting-plane proofs. Discrete Applied Mathematics, 18(1):25–38, 1987.
\bibitem{r5} Stefan S. Dantchev, Nicola Galesi, and Barnaby Martin. Resolution and the binary encoding of combinatorial principles. In Amir Shpilka, editor, 34th Computational Complexity Conference, CCC 2019, July 18-20, 2019, New Brunswick, NJ, USA, volume 137 of LIPIcs, pages 6:1–6:25. Schloss Dagstuhl - Leibniz-Zentrum für Informatik, 2019. doi:10.4230/LIPIcs.CCC.2019.6.
\bibitem{r6} Noah Fleming, Denis Pankratov, Toniann Pitassi, and Robert Robere. Random $\Theta(\log n)$CNFs are hard for cutting planes. In Chris Umans, editor, 58th IEEE Annual Symposium on Foundations of Computer Science, FOCS 2017, Berkeley, CA, USA, October 15-17, 2017, pages 109–120. IEEE Computer Society, 2017. doi:10.1109/FOCS.2017.19.
\bibitem{r7} Ankit Garg, Mika Göös, Pritish Kamath, and Dmitry Sokolov. Monotone circuit lower bounds from resolution. In Ilias Diakonikolas, David Kempe, and Monika Henzinger, editors, Proceedings of the 50th Annual ACM SIGACT Symposium on Theory of Computing, STOC 2018, Los Angeles, CA, USA, June 25-29, 2018, pages 902–911. ACM, 2018. doi:10.1145/ 3188745.3188838.
\bibitem{r8} Michal Garlík and Leszek Aleksander Kolodziejczyk. Some subsystems of constant-depth frege with parity. ACM Trans. Comput. Log., 19(4):29:1–29:34, 2018. doi:10.1145/3243126.
\bibitem{r9} Mika Göös and Toniann Pitassi. Communication lower bounds via critical block sensitivity. In Proceedings of the Forty-Sixth Annual ACM Symposium on Theory of Computing, STOC ’14, page 847–856, New York, NY, USA, 2014. Association for Computing Machinery. doi: 10.1145/2591796.2591838.
\bibitem{r10} Pavel Hrubes and Pavel Pudlák. Random formulas, monotone circuits, and interpolation. In Chris Umans, editor, 58th IEEE Annual Symposium on Foundations of Computer Science, FOCS 2017, Berkeley, CA, USA, October 15-17, 2017, pages 121–131. IEEE Computer Society, 2017. doi:10.1109/FOCS.2017.20.
\bibitem{r11} Pavel Hrubes and Pavel Pudlák. A note on monotone real circuits. Inf. Process. Lett., 131:15–19, 2018. doi:10.1016/j.ipl.2017.11.002.
\bibitem{r12} Trinh Huynh and Jakob Nordström. On the virtue of succinct proofs: amplifying communication complexity hardness to time-space trade-offs in proof complexity. In Howard J. Karloff and Toniann Pitassi, editors, Proceedings of the 44th Symposium on Theory of Computing Conference, STOC 2012, New York, NY, USA, May 19 - 22, 2012, pages 233–248. ACM, 2012. doi:10.1145/2213977.2214000.
\bibitem{r13} Russell Impagliazzo, Toniann Pitassi, and Alasdair Urquhart. Upper and lower bounds for tree-like cutting planes proofs. In Proceedings of the Ninth Annual Symposium on Logic in Computer Science (LICS ’94), Paris, France, July 4-7, 1994, pages 220–228. IEEE Computer Society, 1994. doi:10.1109/LICS.1994.316069.
\bibitem{r14} Dmitry Itsykson and Dmitry Sokolov. Lower bounds for splittings by linear combinations. In Erzsébet Csuhaj-Varjú, Martin Dietzfelbinger, and Zoltán Ésik, editors, Mathematical Foundations of Computer Science 2014 - 39th International Symposium, MFCS 2014, Budapest, Hungary, August 25-29, 2014. Proceedings, Part II, volume 8635 of Lecture Notes in Computer Science, pages 372–383. Springer, 2014. doi:10.1007/978-3-662-44465-8\_32.
\bibitem{r15} Dmitry Itsykson and Dmitry Sokolov. Resolution over linear equations modulo two. Annals of Pure and Applied Logic, 171(1), January 2020. doi:10.1016/j.apal.2019.102722.
\bibitem{r16} Bala Kalyanasundaram and Georg Schnitger. The probabilistic communication complexity of set intersection. SIAM J. Discret. Math., 5(4):545–557, 1992. doi:10.1137/0405044.
\bibitem{r17} Mauricio Karchmer and Avi Wigderson. Monotone circuits for connectivity require superlogarithmic depth. SIAM J. Discret. Math., 3(2):255–265, 1990. doi:10.1137/0403021.
\bibitem{r18} H. Kesten. An introduction to probability theory and its applications, volume i, (william feller). SIAM Review, 11(1):96–96, 1969. doi:10.1137/1011021.
\bibitem{r19} Erfan Khaniki. On proof complexity of resolution over polynomial calculus. Electronic Colloquium on Computational Complexity (ECCC), 27:34, 2020. URL: https://eccc.weizmann. ac.il/report/2020/034.
\bibitem{r20} Jan Krajíček. Interpolation theorems, lower bounds for proof systems, and independence results for bounded arithmetic. J. Symb. Log., 62(2):457–486, 1997. doi:10.2307/2275541.
\bibitem{r21} Jan Krajíček. An exponential lower bound for a constraint propagation proof system based on ordered binary decision diagrams. J. Symb. Log., 73(1):227–237, 2008. doi:10.2178/jsl/ 1208358751.
\bibitem{r22} Jan Krajíček. Randomized feasible interpolation and monotone circuits with a local oracle. J. Mathematical Logic, 18(2):1850012:1–1850012:27, 2018. doi:10.1142/S0219061318500125.
\bibitem{r23} Jan Krajíček. Proof complexity, volume 170. Cambridge University Press, 2019.
\bibitem{r24} Eyal Kushilevitz and Noam Nisan. Communication complexity. Cambridge University Press, 1997.
\bibitem{r25} Edward I Nechiporuk. A boolean function. Engl. transl. in Sov. Phys. Dokl., 10:591–593, 1966.
\bibitem{r26} Vsevolod Oparin. Tight upper bound on splitting by linear combinations for pigeonhole principle. In Nadia Creignou and Daniel Le Berre, editors, Theory and Applications of Satisfiability Testing - SAT 2016 - 19th International Conference, Bordeaux, France, July 5-8, 2016, Proceedings, volume 9710 of Lecture Notes in Computer Science, pages 77–84. Springer, 2016. doi:10.1007/978-3-319-40970-2\_6.
\bibitem{r27} Fedor Part and Iddo Tzameret. Resolution with counting: Dag-like lower bounds and different moduli. In Thomas Vidick, editor, 11th Innovations in Theoretical Computer Science Conference, ITCS 2020, January 12-14, 2020, Seattle, Washington, USA, volume 151 of LIPIcs, pages 19:1–19:37. Schloss Dagstuhl - Leibniz-Zentrum für Informatik, 2020. doi:10.4230/LIPIcs.ITCS.2020.19.
\bibitem{r28} Pavel Pudlák. Lower bounds for resolution and cutting plane proofs and monotone computations. J. Symb. Log., 62(3):981–998, 1997. doi:10.2307/2275583.
\bibitem{r29} Ran Raz and Pierre McKenzie. Separation of the monotone NC hierarchy. Combinatorica, 19(3):403–435, 1999. doi:10.1007/s004930050062.
\bibitem{r30} Ran Raz and Avi Wigderson. Monotone circuits for matching require linear depth. J. ACM, 39(3):736–744, 1992. doi:10.1145/146637.146684.
\bibitem{r31} Alexander A. Sherstov. The multiparty communication complexity of set disjointness. In Proceedings of the forty-fourth annual ACM symposium on Theory of computing, pages 525–548, 2012.
\bibitem{r32} Alexander A. Sherstov. Communication lower bounds using directional derivatives. J. ACM, 61(6), December 2014. doi:10.1145/2629334.
\bibitem{r33} Dmitry Sokolov. Dag-like communication and its applications. In Computer Science Theory and Applications - 12th International Computer Science Symposium in Russia, CSR 2017, Kazan, Russia, June 8-12, 2017, Proceedings, pages 294–307, 2017. doi:10.1007/ 978-3-319-58747-9\_26.
\end{thebibliography}
